% Appendix B.2, input by main.tex.

\FloatBarrier
\subsection{Illustrative walkthrough of the prompt}\label{sec:composition-walkthrough}

The walkthrough follows the sections in Table~\ref{tab:composition-sections} in order. The examples
illustrate the composition of one event in the Part slippage recovery, in which \code{ur5e-3}
picks up the dropped gear. The refusal example comes from a separate composition of the same event.

\subsubsection{Task and Role}\label{sec:comp-task}

The Task and Role section asks the~\gls{llm} to compose only the active recovery event, using the
primitives of the resource that performs it. It explains that an accepted event still requires
a program: the primitive steps must be ordered, their parameters must be compatible, and together
they must reach the complete intended end state. The~\gls{llm} must request any missing facts rather
than guess them.
Listing~\ref{prompt:comp-task} shows the section.

\begin{recbox}{Task and Role}
\begin{lstlisting}
(*@\who{Task and Role}@*)
You are the active replanner authoring primitive_steps.
Current phase: Event-to-Primitive Authoring (event-local).
Author primitive_steps for ONLY the active DES transition, using only the visible primitive catalog for the active resource.
RA primitive_support records independent capability witnesses. Primitive Composition must still establish ordering, compatible parameter bindings, and achievement of the complete expected_end_state. An accepted outline is not an executable program.
The active event is provided inline below. For other facts, emit decision=need_context with refs in context_requests; they are resolved in the next turn under 'Served Context'. Do not guess values. Capability decompositions show how this resource's nominal functions compose primitives. (*@\cut@*)
\end{lstlisting}
\end{recbox}
\promptcaption{The Task and Role section of the primitive composition prompt, the same in every call.}{prompt:comp-task}

\subsubsection{Primitives, Knowledge Base and nominal functions}\label{sec:comp-lists}

The Primitives section lists the names of the primitives of the resource that performs the event.
The program can use only these primitives. As in recovery task synthesis, the Knowledge Base
section lists the facts that the~\gls{llm} can request, such as a primitive's contract, the resource's
model and the geometry of the cell. The Nominal functions section lists the resource's nominal
functions, such as moving home or placing a part. The~\gls{llm} can request the primitive steps used by
each function to see how the resource performs that work in normal production. These sections
list names, and the~\gls{llm} requests further details as needed. This keeps the prompt short.
Listing~\ref{prompt:comp-lists} shows the three sections.

\begin{recbox}{Primitives, Knowledge Base and Nominal functions}
\begin{lstlisting}
(*@\who{Primitives}@*)
[{"name": "compute_pick_targets"}, {"name": "compute_place_targets"}, {"name": "detect_parts"}, {"name": "grasp_part"}, {"name": "move_cartesian"}, {"name": "move_relative"}, {"name": "move_to_named_pose"}, {"name": "release_part"}]

(*@\who{Knowledge Base}@*)
[
{"ref": "/outline_event", "description": "the active outline event (full)"},
{"ref": "/resources/<jid>/snapshot", "description": "full resource snapshot"},
{"ref": "/resources/<jid>/current_pose", "description": "an arm's current tool pose; a KMR's poses are in its snapshot"},
{"ref": "/resources/<jid>/held_part", "description": "currently held part, if any"},
{"ref": "/resources/<jid>/static_capabilities", "description": "as the capture recorded: the places an arm reaches, a machine's process capabilities and configuration; empty for other resources; no poses or coordinates"},
{"ref": "/resources/<jid>/resource_evidence/controller_config/cartesian_motion", "description": "an arm's declared route around its base, as x, y, z points: pick_transit_waypoints and transit_waypoints, and rotation_waypoint, where it turns"},
{"ref": "/parts/<name>/snapshot", "description": "full part state"},
{"ref": "/parts/<name>/observed_pose", "description": "last observed pose for the part, if any"},
{"ref": "/parts/<name>/target", "description": "destination geometry if grounded"},
{"ref": "/primitive_catalog/<jid>", "description": "visible primitive catalog for the resource"},
{"ref": "/primitive_contracts/<name>", "description": "contract card for one visible primitive"},
{"ref": "/projected_snapshot/<jid>", "description": "this resource at the start of the active event"},
{"ref": "/capability_decompositions/<function_name>", "description": "how a nominal function of the active resource composes primitives"},
{"ref": "/regions/<name>", "description": "region bounds in the world frame (or /regions for all)"},
{"ref": "/geometry/<key>", "description": "station poses and sizes, robot and mobile bases, floor extent and obstacles (or /geometry for all)"}
]

(*@\who{Nominal functions}@*)
["move_home", "pick_approach", "pick_grasp", "place_approach", "place_insert"]
\end{lstlisting}
\end{recbox}
\promptcaption{The Primitives, Knowledge Base and Nominal functions sections for \code{ur5e-3}.}{prompt:comp-lists}

\subsubsection{Requested facts}\label{sec:comp-facts}

The next prompt includes the requested facts, which remain available in later calls for the same
event. The resource's model specifies its state variables and nominal events, including the
primitive steps and parameters of those events. A primitive's contract specifies its parameters,
preconditions and effects. A nominal function's decomposition shows the primitive steps used to
perform that work in normal production. Listing~\ref{prompt:comp-facts} shows some of the facts
requested by the~\gls{llm}.

\begin{recbox}{Requested facts}
\begin{lstlisting}
(*@\who{Requested facts}@*)
{"/resources/recovery-resource-3@localhost/recovery_des_model": {"model_type": "extended_finite_automaton", "resource_jid": "recovery-resource-3@localhost", "state_variables": {"held_part": {"scope": "resource", "domain": [null, "KET4_Square_4mm"]}, (*@\cut@*)

{"/primitive_contracts/compute_pick_targets": {"name": "compute_pick_targets", "description": "Compute vertical pick target positions from world object geometry and configured end-effector/TCP feedback without moving. The frames must agree.", (*@\cut@*)
(*@\cut@*)
"/primitive_contracts/grasp_part": {"name": "grasp_part", "description": "Close the gripper and attach the target part as one high-level grasp primitive.", "params": {"model_name": {"type": "string", "x-binding-role": "controller_identifier", "description": "Controller identifier of the part to attach; a semantic part label does not establish this binding"}, "part_name": {"type": "string", "description": "Optional canonical part identifier for held-part tracking."}, "position": {"type": ["number", "null"], "description": "Optional gripper closing position override."}}, "required_params": ["model_name"], "preconditions": {"availability": {"equals": "available"}, "held_part": {"equals": null}, "gripper_state": {"not_equals": "failed"}}, "effects": {"current_state": {"set": "picked"}, "gripper_state": {"set": "closed"}, "held_part": {"set_from_param_any_of": ["part_name", "model_name"]}}, "output_schema": {}},
(*@\cut@*)
"/capability_decompositions/pick_grasp": {"function_name": "pick_grasp", (*@\cut@*) "modeled_transition": "at_pick -> picked", (*@\cut@*) "recovery_visible_steps": [{"primitive": "grasp_part", "params": {"model_name": "<MODEL_NAME_FROM_PART_TARGET>", "part_name": "<PART>"}}, (*@\cut@*)
(*@\cut@*)
"/parts/gear_small/observed_pose": {"x": 0.24992089458069847, "y": 0.44999843654915983, "z": 1.0249998633209128, "qx": -2.4896830184808793e-05, "qy": -1.7065517433252323e-06, "qz": 0.0006494841124204398, "qw": 0.9999997887737893}}
\end{lstlisting}
\end{recbox}
\promptcaption{Some of the requested facts in the second call for the gear pick.}{prompt:comp-facts}

\subsubsection{Validation feedback}\label{sec:comp-feedback}

If the previous program was refused, the Validation feedback section appears before the active
event and is marked to be read first. It identifies the event and provides a refusal code and a
reason, specifying which step was refused and why. Every resource's program is checked step by
step against the parameters, preconditions and effects of its primitives. For a robot, MoveIt also
plans the motions using the model and scene recorded at the failure, without moving the robot.
The~\gls{llm} uses this feedback to revise and resubmit the program. Listing~\ref{prompt:comp-feedback}
shows the feedback on a refused program.

\begin{recbox}{Validation feedback}
\begin{lstlisting}
(*@\who{Validation feedback}@*)
- RECOVERY_SEQ3 / recovery-resource-3@localhost [composition_infeasible]: move_to_named_pose step 0 (ur5e-3): MoveIt's Cartesian path on the captured model stopped 30% of the way along waypoint 1 of 1 toward (0.4010, 0.2023, 1.5011); the last reached tool pose is (-0.1010, 0.4096, 1.5010), 0.136 m from ur5e-3's base axis; the IK search found no joint solution there (error code -6). Native reason: Cartesian path is incomplete; no complete motion evidence is available
\end{lstlisting}
\end{recbox}
\promptcaption{Validation feedback on a refused program for the gear pick, from a separate composition of the same event.}{prompt:comp-feedback}

\subsubsection{Active event and predefined configurations}\label{sec:comp-event}

The Active event section presents the recovery event as accepted in recovery task synthesis.
It identifies the event, its resource and its part, then specifies the starting state, intended
end state and the~\gls{llm}'s rationale. The Predefined configurations section follows,
listing the configurations declared by the resource that a step can reference by name.
Listing~\ref{prompt:comp-event} shows both sections for the gear pick.

\begin{recbox}{Active event and Predefined configurations}
\begin{lstlisting}
(*@\who{Active event}@*)
{"outline_id": "RECOVERY_SEQ3", "resource_jid": "recovery-resource-3@localhost", "action": "pick_dropped_gear_small_ur5e3", "part_name": "gear_small", "start_state_id": null, "end_state_id": null}

{"outline_id": "RECOVERY_SEQ3", "resource_jid": "recovery-resource-3@localhost", "event_name": "pick_dropped_gear_small_ur5e3", "description": "gear_small's measured position is 0.255 m from ur5e-3's base in xy, at z 1.025, which is inside its envelope (radius 0.886 m, z 0.971(*@\textendash@*)1.121); the pick target comes from the part's recorded position, so no pose is given.", "part_name": "gear_small", "target_ref": "recovery-resource-3@localhost",
 "expected_start_state": {"held_part": null, "part_location": null, "resource_location": "Buffer For Machined parts", "resource_state": "placed"},
 "expected_end_state": {"resource_state": "picked", "held_part": "gear_small", "resource_location": null, "part_location": "recovery-resource-3@localhost"}}
(*@\cut@*)
(*@\who{Predefined configurations}@*)
["home"]
- move_to_named_pose may only use named poses advertised for the active resource.
\end{lstlisting}
\end{recbox}
\promptcaption{The Active event and Predefined configurations sections for the gear pick.}{prompt:comp-event}

\subsubsection{Composition Rules}\label{sec:comp-rules}

The Composition Rules section specifies the requirements for the reply and for an accepted
program. The program can use only the listed primitives, and every number must be expressed as
a reference to a requested fact. Every intended end-state value must be established by the
declared effects of the steps or by what the checks find. Listing~\ref{prompt:comp-rules} shows the section.

\begin{recbox}{Composition Rules}
\begin{lstlisting}
(*@\who{Composition Rules}@*)
- Keep `plan` to at most two short sentences and `rationale` to one short sentence.
- Author primitive_steps for ONLY the active DES transition shown above.
- outline_id and resource_jid MUST exactly match the active transition.
- primitive_steps[*].primitive MUST be a name in Visible Primitive Catalog.
- Each step MUST include params. Every number in params is a {"context_ref": ...}; strings and booleans may be literals from served context.
- compute_pick_targets publishes event_facts.pick_targets.<part>; compute_place_targets publishes event_facts.place_targets.<part>.
- need_context: refs in context_requests, primitive_steps empty.
- primitive_steps_ready: non-empty primitive_steps. need_primitive_revision: say why in rationale.

RA composition validation: author an ordered program that establishes every intended resource and part postcondition from declared effects or what the checks find. (*@\cut@*)
\end{lstlisting}
\end{recbox}
\promptcaption{The Composition Rules section, the same in every call apart from the resource it names.}{prompt:comp-rules}

\subsubsection{System message}\label{sec:comp-system}

The system message requires a single JSON object and specifies its schema. Every reply contains
a short plan, the event identifiers and a decision. The decision indicates whether the~\gls{llm} is
requesting facts, providing a program of primitive steps, or reporting that the event cannot be
composed from the primitives. In the last case, the reply includes the reason. Each step names a
primitive and its parameters. Listing~\ref{prompt:comp-system} shows the system message.

\begin{recbox}{System message}
\begin{lstlisting}
Reply with only one JSON object that follows this JSON schema:
{"type": "object", "properties": {"plan": {"type": "string"}, "decision": {"type": "string", "enum": ["primitive_steps_ready", "need_context", "need_primitive_revision"]}, "outline_id": {"type": "string"}, "resource_jid": {"type": "string"}, "context_requests": {"type": "array", "items": {"type": "string"}}, "rationale": {"type": "string"}, "notes": {"type": "array", "items": {"type": "string"}}, "primitive_steps": {"type": "array", "items": {"type": "object", "properties": {"primitive": {"type": "string"}, "params": {"type": "object"}}, "required": ["primitive", "params"]}}}, "required": ["plan", "decision", "outline_id", "resource_jid"]}
\end{lstlisting}
\end{recbox}
\promptcaption{The system message of the primitive composition prompt, the same in every call.}{prompt:comp-system}

\subsubsection{Replies of the~\texorpdfstring{\gls{llm}}{LLM}}\label{sec:comp-replies}

Listing~\ref{reply:comp-examples} shows \gls{llm} replies from two calls of the gear pick composition.
In the first call, the~\gls{llm} requests facts. In the third call, it gives the program that the~\gls{ra}
accepted.

\begin{recbox}{Replies of the~\gls{llm}}
\begin{lstlisting}
(*@\who{Call 1}@*)
{"plan": "Gather the primitive contracts, pick decompositions, and part/resource state needed to compose a pick of gear_small by ur5e-3. Then author the ordered steps: compute targets, approach, grasp, retreat.", "decision": "need_context", (*@\cut@*) "context_requests": ["/primitive_contracts/compute_pick_targets", "/primitive_contracts/move_cartesian", "/primitive_contracts/grasp_part", "/primitive_contracts/move_relative", "/primitive_contracts/move_to_named_pose", "/capability_decompositions/pick_approach", "/capability_decompositions/pick_grasp", "/projected_snapshot/recovery-resource-3@localhost", "/parts/gear_small/snapshot", "/parts/gear_small/observed_pose", "/resources/recovery-resource-3@localhost/recovery_des_model"], (*@\cut@*)}
(*@\who{Call 3}@*)
{"plan": "Compute gear_small's pick targets from its recorded observation, then approach, descend and grasp it. Finish by retreating to the computed retreat pose.", "decision": "primitive_steps_ready", (*@\cut@*) "primitive_steps": [
{"primitive": "compute_pick_targets", "params": {"part_name": "gear_small"}},
{"primitive": "move_cartesian", "params": {"x": {"context_ref": "event_facts.pick_targets.gear_small.approach_pose.x"}, (*@\cut@*)}},
{"primitive": "move_cartesian", "params": {"x": {"context_ref": "event_facts.pick_targets.gear_small.target_pose.x"}, (*@\cut@*)}},
{"primitive": "grasp_part", "params": {"model_name": {"context_ref": "/parts/gear_small/target/model_name"}, "part_name": "gear_small"}},
{"primitive": "move_cartesian", "params": {"x": {"context_ref": "event_facts.pick_targets.gear_small.retreat_pose.x"}, (*@\cut@*)}}]}
\end{lstlisting}
\end{recbox}
\promptcaption{Two \gls{llm} replies for the gear pick. The program given in the second call was refused.}{reply:comp-examples}
